\documentclass[12pt,oneside]{report}

% ── Page geometry ──
\usepackage[letterpaper,left=1.5in,right=1in,top=1in,bottom=1in,includefoot]{geometry}

% ── Typography ──
\usepackage[T1]{fontenc}
\usepackage{lmodern}
\usepackage[utf8]{inputenc}
\usepackage{microtype}

% ── Colors ──
\usepackage[table]{xcolor}
\definecolor{jhublue}{HTML}{002D72}
\definecolor{jhugold}{HTML}{CFB87C}
\definecolor{jhuslate}{HTML}{2D2A26}
\definecolor{jhulight}{HTML}{F2EFE9}

% ── Graphics and boxes ──
\usepackage{graphicx}
\usepackage{tcolorbox}
\tcbuselibrary{skins,breakable}
\usepackage{tikz}
\usetikzlibrary{positioning,calc,shapes.geometric,backgrounds,patterns}

% ── Math and symbols ──
\usepackage{amsmath,amssymb,amsthm}

% ── Tables ──
\usepackage{booktabs}
\usepackage{array}
\usepackage{tabularx}

% ── Lists ──
\usepackage{enumitem}

% ── Headers and footers ──
\usepackage{fancyhdr}
\usepackage{emptypage}

% ── Cross-references and hyperlinks ──
\usepackage[hidelinks,colorlinks=false]{hyperref}

% ── Figure and table captions ──
\usepackage[margin=10pt,font={sf,footnotesize},labelfont={bf,color=jhublue}]{caption}

% ── Bibliography ──
\usepackage[round]{natbib}

% ── Epigraph ──
\usepackage{epigraph}

% ── Multi-column ──
\usepackage{multicol}

% ── Line spacing ──
\usepackage{setspace}
\setstretch{1.1}

% ── Thesis / frontmatter ──
\usepackage{titling}

% ── Headers and footers ──
\pagestyle{fancy}
\fancyhf{}
\fancyhead[L]{\footnotesize\color{jhublue}\nouppercase{\leftmark}}
\fancyhead[R]{\footnotesize\color{jhublue}\thepage}
\renewcommand{\headrulewidth}{0.4pt}
\renewcommand{\headrule}{\hbox to\headwidth{\color{jhublue}\leaders\hrule height \headrulewidth\hfill}}
\fancypagestyle{plain}{
  \fancyhf{}
  \fancyhead[R]{\footnotesize\color{jhublue}\thepage}
  \renewcommand{\headrulewidth}{0pt}
}

% ── Chapter formatting ──
\usepackage{titlesec}
\titleformat{\chapter}[display]
  {\normalfont\centering\color{jhublue}}
  {\scshape\Large C\,H\,A\,P\,T\,E\,R\\[4pt]\Huge\thechapter}
  {12pt}
  {\Huge\bfseries}
\titlespacing*{\chapter}{0pt}{40pt}{36pt}

\titleformat{\section}
  {\normalfont\Large\bfseries\color{jhublue}}
  {\thesection}
  {1em}
  {}
\titlespacing*{\section}{0pt}{24pt}{10pt}

\titleformat{\subsection}
  {\normalfont\large\bfseries\color{jhuslate}}
  {\thesubsection}
  {1em}
  {}
\titlespacing*{\subsection}{0pt}{18pt}{8pt}

% ── Theorem environments ──
\newtheorem{theorem}{Theorem}[chapter]
\newtheorem{lemma}[theorem]{Lemma}
\newtheorem{corollary}[theorem]{Corollary}
\newtheorem{proposition}[theorem]{Proposition}
\theoremstyle{definition}
\newtheorem{definition}[theorem]{Definition}
\newtheorem{example}[theorem]{Example}
\theoremstyle{remark}
\newtheorem{remark}[theorem]{Remark}

% ── TikZ block style ──
\tikzset{
  block/.style = {rectangle, draw=jhublue, thick, fill=jhulight,
    text width=8em, text centered, rounded corners, minimum height=3em},
  line/.style = {draw=jhublue, thick, ->, >=stealth},
}

% ── epigraph setup ──
\setlength{\epigraphwidth}{0.6\textwidth}
\renewcommand{\epigraphrule}{0pt}
\renewcommand{\textflush}{flushepinormal}

% ======================================================================
\begin{document}

% ── Half-title ──
\thispagestyle{empty}
\vspace*{0.4\textheight}
\begin{center}
  {\Huge\bfseries\color{jhublue} Decentralized Coordination}\\[14pt]
  {\Huge\bfseries\color{jhublue} for Multi-Robot Systems}\\[14pt]
  {\Huge\bfseries\color{jhublue} in Uncertain Environments}
\end{center}
\clearpage

% ── TITLE PAGE ──
\thispagestyle{empty}
\begin{center}

\vspace*{18pt}

{\Huge\bfseries\color{jhublue} Decentralized Coordination}\\[10pt]
{\Huge\bfseries\color{jhublue} for Multi-Robot Systems}\\[10pt]
{\Huge\bfseries\color{jhublue} in Uncertain Environments}

\vspace{32pt}

{\large\color{jhuslate} by}

\vspace{16pt}

{\Large\bfseries\color{jhuslate} Margaret E. Calloway}

\vspace{60pt}

{
\color{jhublue}\hrule height 0.5pt
\vspace{16pt}
}

{\normalsize\color{jhuslate}
A dissertation submitted to Johns Hopkins University\\
in conformity with the requirements for the degree of\\
Doctor of Philosophy}

\vspace{16pt}

{\normalsize\color{jhuslate}
Baltimore, Maryland\\
May 2026}

\vspace{40pt}

{
\color{jhublue}\hrule height 0.5pt
}

\vspace{24pt}

{\footnotesize\color{jhublue}
\textcopyright\ 2026 by Margaret E. Calloway\\
All rights reserved}

\vfill

{\scriptsize\color{gray}
Unofficial template --- not affiliated with or endorsed by the university.\\
Sample content only.}

\end{center}
\clearpage

% ── ABSTRACT ──
\thispagestyle{empty}
\phantomsection
\addcontentsline{toc}{chapter}{Abstract}
\begin{center}
  {\LARGE\bfseries\color{jhublue} Abstract}
\end{center}
\vspace{16pt}

Multi-robot systems operating in uncertain, dynamic environments demand coordination
strategies that are both robust to failure and scalable with team size. This dissertation
presents a novel framework for decentralized multi-robot coordination that combines
distributed consensus algorithms with learned uncertainty models to achieve reliable
collective behavior without centralized control. We formalize the problem as a
decentralized Markov decision process augmented with belief-state tracking,
and we derive convergence guarantees for a class of consensus protocols operating
over time-varying communication graphs. The principal contributions are threefold.
\textit{First}, we introduce the \textbf{B}elief-\textbf{A}ugmented \textbf{D}ecentralized
\textbf{C}oordination (BADC) framework, which enables robots to maintain local
estimates of both environmental state and teammate intent while tolerating intermittent
communication. \textit{Second}, we develop a polynomial-time planning algorithm that
exploits submodular structure in the team objective to produce near-optimal joint
actions with provable approximation bounds. \textit{Third}, we validate the framework
through extensive simulation studies across six benchmark domains --- including
search-and-rescue, formation control, and cooperative transport --- and through physical
experiments with a fleet of eight ground robots navigating a cluttered warehouse
environment. Results demonstrate that BADC reduces coordination latency by
$42\%$ relative to the nearest centralized baseline while matching its task-completion
rate, and that the system degrades gracefully under communication dropout rates
of up to $35\%$. The implications of this work extend beyond robotics to any
distributed system where agents must coordinate under partial observability,
including autonomous vehicle fleets, drone swarms, and edge-computing clusters.

\vspace{28pt}

\noindent{\color{jhublue}\textbf{Primary Reader:}} \hspace{4pt} Dr.\ Alistair V. Chen\\
\noindent{\color{jhublue}\textbf{Secondary Reader:}} \hspace{4pt} Dr.\ Priya N. Okonkwo

\clearpage

% ── ACKNOWLEDGMENTS ──
\thispagestyle{empty}
\phantomsection
\addcontentsline{toc}{chapter}{Acknowledgments}
\begin{center}
  {\LARGE\bfseries\color{jhublue} Acknowledgments}
\end{center}
\vspace{16pt}

I am deeply grateful to my advisor, Dr.\ Alistair Chen, whose patience, insight,
and steady encouragement shaped every stage of this work. His belief that good theory
must answer to real hardware kept this project grounded and gave it purpose.

I thank the members of my dissertation committee --- Dr.\ Priya Okonkwo, Dr.\ Samuel
Becker, and Dr.\ Helena Vargas --- for their careful reading and probing questions.
Their feedback sharpened the presentation substantially.

The Laboratory for Autonomous Systems provided both the physical and intellectual
home for this research. I am especially indebted to Karim Nassar and Dr.\ Lina
Thorvaldsen for countless hours of debugging, whiteboard sessions, and late-night
simulation runs. This work benefited from financial support through the Vertex
Foundation Graduate Fellowship, the Nimbus Institute for Robotics grant \#NB-2024-1139,
and a Synapse Technologies research gift.

Finally, I thank my parents, Elena and Thomas Calloway, and my brother Julian,
whose quiet confidence in me never wavered. To Mira --- you made Baltimore feel
like home.

\clearpage

% ── TABLE OF CONTENTS ──
\phantomsection
\addcontentsline{toc}{chapter}{Table of Contents}
\renewcommand{\contentsname}{\centering\LARGE\bfseries\color{jhublue} Table of Contents}
\tableofcontents
\cleardoublepage

% ── LIST OF FIGURES ──
\phantomsection
\addcontentsline{toc}{chapter}{List of Figures}
\renewcommand{\listfigurename}{\centering\LARGE\bfseries\color{jhublue} List of Figures}
\listoffigures
\addcontentsline{toc}{chapter}{List of Tables}
\renewcommand{\listtablename}{\centering\LARGE\bfseries\color{jhublue} List of Tables}
\listoftables
\cleardoublepage

% ======================================================================
% CHAPTER 1: INTRODUCTION
% ======================================================================
\chapter{Introduction}
\pagenumbering{arabic}
\setcounter{page}{1}

\epigraph{\itshape The whole is greater than the sum of its parts.}{--- Aristotle, \textit{Metaphysics}}

\section{Motivation}

Teams of autonomous robots are poised to transform domains ranging from disaster
response to planetary exploration. A swarm of small, inexpensive robots can search
a collapsed building faster than a single sophisticated platform; a fleet of autonomous
underwater vehicles can map an ocean basin with greater coverage and resilience
than any individual sensor package. Yet realizing these benefits depends on solving
a fundamental challenge: how should a collection of independent agents coordinate
their actions when each agent has only a partial, noisy view of the world and
cannot reliably communicate with its teammates?

Centralized approaches --- where a single planner computes actions for all agents ---
offer theoretical simplicity and, in some cases, optimality guarantees. However,
they suffer from three well-known weaknesses: they introduce a single point of
failure, they scale poorly with team size, and they assume a communication
infrastructure that is unavailable in many practical settings. The alternative,
decentralized coordination, distributes decision-making across agents and promises
fault tolerance, scalability, and communication efficiency. The cost is a more
difficult planning problem, because each agent must reason not only about the
environment but also about the beliefs, intentions, and likely future actions of
its teammates.

This dissertation addresses that cost. We develop a framework that makes
decentralized coordination both principled and practical for teams of robots
operating in uncertain, dynamic environments.

\section{Problem Statement}

Consider a team of $n$ robots indexed by $i \in \{1, \ldots, n\}$, each operating
in a shared environment modeled as a Markov decision process (MDP). Unlike the
standard single-agent MDP, however, the transition dynamics depend on the joint
action of all agents, and each robot receives only a local, partial observation
of the global state. Formally, the system is a \textbf{Dec-POMDP} --- a
decentralized partially observable Markov decision process --- defined by the tuple
$\langle I, S, \{A_i\}, T, \{O_i\}, \Omega, R, \gamma \rangle$, where:

\vspace{6pt}
\begin{tcolorbox}[
  colback=jhulight,
  colframe=jhublue,
  arc=4pt,
  boxrule=0.8pt,
  left=16pt,
  right=16pt,
  top=10pt,
  bottom=10pt
]
\vspace{-4pt}
\begin{itemize}[leftmargin=*,itemsep=3pt,topsep=0pt]
  \item $I = \{1, \ldots, n\}$ is the set of agents.
  \item $S$ is a finite set of world states.
  \item $A_i$ is the action set of agent $i$; the joint action space is $\mathcal{A} = \times_i A_i$.
  \item $T(s' \mid s, \mathbf{a})$ is the transition probability from state $s$ to $s'$ under joint action $\mathbf{a}$.
  \item $O_i$ is the observation set of agent $i$.
  \item $\Omega(\mathbf{o} \mid s, \mathbf{a})$ is the probability of joint observation $\mathbf{o}$.
  \item $R(s, \mathbf{a})$ is the immediate team reward.
  \item $\gamma \in [0, 1)$ is the discount factor.
\end{itemize}
\vspace{-4pt}
\end{tcolorbox}
\vspace{6pt}

The objective is to find a tuple of local policies $\pi = (\pi_1, \ldots, \pi_n)$
that maximizes the expected discounted sum of rewards:
\begin{equation}
  J(\pi) = \mathbb{E}\left[ \sum_{t=0}^{\infty} \gamma^t \, R(s_t, \mathbf{a}_t) \;\middle|\; \pi \right].
  \label{eq:objective}
\end{equation}

Solving the Dec-POMDP exactly is NEXP-complete in the worst case
\citep{bernstein2002complexity}, which motivates approximate, structured approaches.

\section{Contributions}

This dissertation makes the following contributions:

\begin{enumerate}[leftmargin=*,itemsep=6pt,topsep=6pt]
  \item \textbf{The BADC framework} (Chapter~\ref{ch:badc}): A belief-augmented
    architecture in which each robot maintains a factored representation of
    environmental uncertainty and teammate intent. We prove that, under mild
    connectivity assumptions, the local beliefs converge to a neighborhood
    of the centralized Bayesian posterior.

  \item \textbf{Submodular planning algorithm} (Chapter~\ref{ch:planning}):
    We exploit latent submodularity in the team objective to construct a greedy
    planner with approximation ratio $(1 - 1/e)$ and runtime $O(n^2 k \log k)$,
    where $k$ is the planning horizon. This represents an exponential improvement
    over the naive joint-action enumeration.

  \item \textbf{Empirical evaluation} (Chapter~\ref{ch:evaluation}):
    Comprehensive experiments across six simulation benchmarks and one physical
    deployment show consistent improvements in coordination latency,
    communication efficiency, and graceful degradation under failure.
\end{enumerate}

\section{Thesis Outline}

The remainder of this dissertation is organized as follows. Chapter~\ref{ch:related}
surveys related work in multi-agent planning, distributed consensus, and robot
coordination. Chapter~\ref{ch:badc} presents the BADC framework and establishes its
theoretical properties. Chapter~\ref{ch:planning} develops the submodular planning
algorithm. Chapter~\ref{ch:evaluation} describes the experimental methodology and
presents results. Chapter~\ref{ch:conclusion} summarizes findings and discusses
directions for future work.

% ======================================================================
% CHAPTER 2: RELATED WORK
% ======================================================================
\chapter{Related Work}
\label{ch:related}

This chapter situates the present work within three bodies of literature:
multi-agent decision-making under uncertainty, distributed consensus protocols,
and multi-robot coordination architectures.

\section{Multi-Agent Decision-Making}

The study of decentralized decision-making dates to the earliest formulations
of the Dec-POMDP \citep{bernstein2002complexity}. Exact solution methods,
including dynamic programming for finite horizons \citep{hansen2004dynamic}
and point-based value iteration \citep{spaan2005perseus}, are limited to
small teams ($n \leq 3$) and small state spaces. Approximate methods fall
broadly into two categories.

\textbf{Policy-search approaches} optimize the parameters of a restricted
policy class using reinforcement learning or evolutionary strategies.
Recent work on multi-agent actor-critic methods \citep{lowe2017maddpg} and
counterfactual multi-agent policy gradients \citep{foerster2018counterfactual}
has demonstrated impressive results in simulated domains, but these methods
typically assume centralized training with access to global state --- an
assumption relaxed in our work.

\textbf{Structured decomposition approaches} exploit independence or
interaction sparsity to factor the joint value function. Transition-independent
Dec-MDPs \citep{becker2004transition} and ND-POMDPs
\citep{nair2005networked} achieve tractability by restricting agent
interactions to a sparse graph, an idea we extend in the BADC framework
(Section~\ref{sec:badc-structure}).

\section{Distributed Consensus}

Consensus algorithms enable a group of agents to agree on shared quantities
using only local communication. The foundational results of
\citet{olfati2007consensus} establish convergence conditions for
time-invariant communication topologies, while \citet{ren2005consensus}
extend these to switching graphs. Our work builds on the push-sum protocol
\citep{kempe2003gossip} and its robust variants \citep{hadjicostis2017average},
which we adapt to the belief-tracking context of BADC.

\section{Multi-Robot Coordination Architectures}

Three architectural paradigms dominate the multi-robot coordination literature.

\begin{table}[htbp]
\centering
\caption{Comparison of multi-robot coordination paradigms.}
\label{tab:paradigms}
\vspace{4pt}
\begin{tabular}{@{}p{2.5cm} p{3.2cm} p{3.2cm} p{3.2cm}@{}}
\toprule
\color{jhublue}\textbf{Criterion} &
\color{jhublue}\textbf{Centralized} &
\color{jhublue}\textbf{Decentralized} &
\color{jhublue}\textbf{BADC (this work)} \\
\midrule
Fault tolerance    & Single point of failure    & Robust                     & Robust                        \\
Scalability        & $O(n^2)$ communication      & $O(n)$ communication        & $O(n \log n)$ communication   \\
Optimality guarantee& Optimal given full state    & Suboptimal in general       & $(1-1/e)$-approximation       \\
Comm.\ assumption  & Full connectivity           & Sparse, intermittent        & Sparse, intermittent          \\
Planning runtime   & Exponential in $n$          & Heuristic-dependent         & $O(n^2 k \log k)$             \\
\bottomrule
\end{tabular}
\vspace{4pt}
\end{table}

Table~\ref{tab:paradigms} summarizes the trade-offs among these paradigms.
The BADC framework occupies a middle ground, combining the robustness and
scalability of decentralized methods with approximation guarantees that
approach centralized optimality under favorable conditions.

% ======================================================================
% CHAPTER 3: BADC FRAMEWORK
% ======================================================================
\chapter{The BADC Framework}
\label{ch:badc}

\section{Overview}

The Belief-Augmented Decentralized Coordination (BADC) framework consists of
three interacting components: a \textit{local belief tracker}, a
\textit{consensus module}, and a \textit{planning module}. Figure~\ref{fig:badc-arch}
illustrates the data flow among these components.

\begin{figure}[htbp]
\centering
\resizebox{\linewidth}{!}{%
\begin{tikzpicture}[node distance=2.2cm, auto]
  \node [block] (obs) {Local Observations};
  \node [block, right of=obs, xshift=3.5cm] (belief) {Belief Tracker};
  \node [block, right of=belief, xshift=3.5cm] (consensus) {Consensus Module};
  \node [block, below of=obs, yshift=-1.8cm] (action) {Action Selection};
  \node [block, below of=belief, yshift=-1.8cm] (planner) {Submodular Planner};
  \node [block, below of=consensus, yshift=-1.8cm] (execute) {Execution};

  \draw[->, line] (obs) -- (belief);
  \draw[->, line] (belief) -- node[above] {$b_i^t$} (consensus);
  \draw[->, line] (consensus.east) -- ++(1cm,0) |- node[right,yshift=6pt] {to peers} (consensus.east);
  \draw[->, line] (belief) -- (planner);
  \draw[->, line] (planner) -- (action);
  \draw[->, line] (action) -- (execute);

  \draw[->, line, dashed] (consensus.south) -- ++(0,-0.8cm) -| node[below right,yshift=-4pt] {$\tilde{b}_i^t$} (planner.north);
  \draw[->, line, dashed] (execute.west) -- ++(-5cm,0) |- (obs.south);
\end{tikzpicture}%
}
\caption{Architecture of the BADC framework. Solid lines indicate feedforward
data flow; dashed lines indicate feedback and inter-agent communication.}
\label{fig:badc-arch}
\end{figure}

\section{Belief Tracking with Structured Uncertainty}
\label{sec:badc-structure}

Each agent $i$ maintains a factored belief state $b_i^t = (b_i^{\text{env}}, b_i^{\text{team}})$,
where $b_i^{\text{env}} \in \Delta(S)$ tracks the distribution over world states
and $b_i^{\text{team}}$ tracks the estimated intent of each teammate. The
separation is motivated by the observation that environmental uncertainty updates
are driven by local sensor measurements, while teammate-intent uncertainty is
resolved primarily through communication.

The environmental belief update follows a standard Bayesian filter:
\begin{equation}
  b_i^{\text{env}}(s') = \eta \; \Omega(o_i \mid s', a_i) \sum_{s \in S} T(s' \mid s, \mathbf{\hat{a}}) \; b_i^{\text{env}}(s),
  \label{eq:bayes}
\end{equation}
where $\mathbf{\hat{a}}$ denotes agent $i$'s estimate of the joint action
(derived from $b_i^{\text{team}}$) and $\eta$ is a normalization constant.

\begin{definition}[Consensus Belief]
  Let $\mathcal{G}^t = (V, E^t)$ be the communication graph at time $t$, where
  an edge $(i, j) \in E^t$ indicates that agents $i$ and $j$ can exchange messages.
  The \emph{consensus belief} $\tilde{b}_i^t$ is obtained by applying the push-sum
  protocol to the local beliefs:
  \begin{equation}
    \tilde{b}_i^{t}(x) = \frac{\sum_{j \in \mathcal{N}_i^t \cup \{i\}} w_{ij}^t \, b_j^t(x)}
                               {\sum_{j \in \mathcal{N}_i^t \cup \{i\}} w_{ij}^t},
  \end{equation}
  where $\mathcal{N}_i^t = \{j \mid (i, j) \in E^t\}$ and $w_{ij}^t$ are
  Metropolis weights based on node degrees.
\end{definition}

\begin{theorem}[Convergence of BADC Beliefs]
  \label{thm:convergence}
  Assume the communication graph sequence $\{\mathcal{G}^t\}_{t \geq 0}$ is
  uniformly jointly connected. Then for any $\varepsilon > 0$, there exists
  $T_0$ such that for all $t \geq T_0$ and all agents $i, j$:
  \begin{equation}
    \Vert \tilde{b}_i^t - \tilde{b}_j^t \Vert_1 \leq \varepsilon.
  \end{equation}
  Moreover, the rate of convergence is geometric in the number of consensus rounds
  per time step.
\end{theorem}

\begin{proof}[Proof sketch]
  Under the uniformly jointly connected assumption, the product of stochastic
  matrices $W^t W^{t-1} \cdots W^{t-B+1}$ is scrambling for sufficiently large
  window $B$. The result follows from the coefficient of ergodicity bound
  \citep{seneta2006nonnegative} applied to the push-sum dynamics. A complete
  proof appears in Appendix~\ref{app:convergence}.
\end{proof}

\section{Communication Efficiency}

A key design choice in BADC is \emph{event-triggered} communication: agents
broadcast their local beliefs only when the KL-divergence between the current
belief and the last transmitted belief exceeds a threshold $\delta$:

\begin{equation}
  D_{\mathrm{KL}}(b_i^t \;\Vert\; b_i^{\text{last}}) \geq \delta.
  \label{eq:trigger}
\end{equation}

This mechanism reduces communication volume by $60\%$--$80\%$ in typical
scenarios while introducing negligible coordination error (see
Chapter~\ref{ch:evaluation}).

% ======================================================================
% CHAPTER 4: SUBMODULAR PLANNING
% ======================================================================
\chapter{Submodular Planning for Team Coordination}
\label{ch:planning}

\section{Submodularity in Team Objectives}

A set function $f: 2^V \to \mathbb{R}$ is \emph{submodular} if for all
$A \subseteq B \subseteq V$ and $x \notin B$, it satisfies diminishing returns:
\begin{equation}
  f(A \cup \{x\}) - f(A) \geq f(B \cup \{x\}) - f(B).
  \label{eq:submodular}
\end{equation}

Many natural team objectives exhibit this property. Consider a coverage task
where the reward is the number of grid cells visited by at least one robot.
Adding a robot to a small team increases coverage substantially; adding the
same robot to an already-large team yields a smaller marginal gain.

\begin{proposition}[Submodularity of Coverage Reward]
  Let $R(\mathbf{a}) = |\bigcup_{i=1}^n \text{Coverage}(a_i)|$ be the coverage
  reward. Then $R$ is monotone submodular in the set of activated robots.
\end{proposition}

\begin{proof}
  Coverage is the cardinality of a union of sets, and the cardinality-of-union
  function is known to be submodular for any collection of subsets
  \citep{krause2014submodularity}. Monotonicity follows because adding robots
  can only increase (or leave unchanged) the covered area.
\end{proof}

\section{The Greedy-BADC Planner}

Algorithm~\ref{alg:greedy} presents the greedy planning procedure. At each
decision epoch, each agent independently computes a candidate action by
greedily maximizing the marginal contribution to the estimated team objective,
using the consensus belief $\tilde{b}_i^t$.

\begin{tcolorbox}[
  colback=jhulight,
  colframe=jhublue,
  arc=4pt,
  boxrule=0.8pt,
  left=16pt,
  right=16pt,
  top=10pt,
  bottom=10pt,
  title={\color{jhublue}\textbf{Algorithm 1: Greedy-BADC Planner}},
  fonttitle=\bfseries
]
\vspace{2pt}
\textbf{Input:} Consensus belief $\tilde{b}_i^t$, action space $A_i$,
  objective $F$\\
\textbf{Output:} Selected action $a_i^*$
\vspace{4pt}
\begin{enumerate}[leftmargin=*,itemsep=2pt,topsep=2pt]
  \item Initialize candidate set $\mathcal{C} \leftarrow \varnothing$
  \item \textbf{for} each teammate $j \neq i$ \textbf{do}\\[-2pt]
        \hspace*{12pt} $\hat{a}_j \leftarrow \arg\max_{a \in A_j} F(\tilde{b}_i^t, \mathcal{C} \cup \{a\})$\\[-2pt]
        \hspace*{12pt} $\mathcal{C} \leftarrow \mathcal{C} \cup \{\hat{a}_j\}$
  \item \textbf{end for}
  \item $a_i^* \leftarrow \arg\max_{a \in A_i} F(\tilde{b}_i^t, \mathcal{C} \cup \{a\})$
  \item \textbf{return} $a_i^*$
\end{enumerate}
\vspace{2pt}
\end{tcolorbox}
\vspace{8pt}

\begin{theorem}[Approximation Guarantee]
  \label{thm:approx}
  Let $F$ be a monotone submodular set function and let $\pi^*$ be an optimal
  joint policy. The Greedy-BADC planner achieves expected reward
  $J(\pi_{\text{BADC}}) \geq (1 - 1/e) \, J(\pi^*)$ under the consensus belief
  model, provided that the communication graph satisfies the conditions of
  Theorem~\ref{thm:convergence}.
\end{theorem}

\begin{proof}
  The result follows from the classical guarantee for greedy maximization of
  monotone submodular functions \citep{nemhauser1978analysis} combined with
  the belief-convergence bound of Theorem~\ref{thm:convergence}. The
  approximation error introduced by imperfect consensus is bounded by
  $O(\varepsilon)$ and can be made arbitrarily small by increasing the
  number of consensus rounds.
\end{proof}

\section{Complexity Analysis}

The per-agent runtime of Algorithm~\ref{alg:greedy} is $O(n \cdot |A_{\max}| \cdot T_F)$,
where $|A_{\max}| = \max_i |A_i|$ and $T_F$ is the cost of evaluating the
objective $F$. For the coverage objective, $T_F = O(m \log m)$ using a
union-find data structure over $m$ grid cells, yielding an overall complexity
of $O(n \cdot |A_{\max}| \cdot m \log m)$ per decision epoch. This compares
favorably to the $O(|A_{\max}|^n \cdot m)$ cost of exhaustive joint-action search.

% ======================================================================
% CHAPTER 5: EXPERIMENTAL EVALUATION
% ======================================================================
\chapter{Experimental Evaluation}
\label{ch:evaluation}

\section{Simulation Benchmarks}

We evaluate BADC against four baselines --- centralized MDP (C-MDP), independent
Q-learning (IQL) \citep{tan1993multi}, consensus-based auction (CBBA)
\citep{choi2009consensus}, and decentralized MCTS (D-MCTS)
\citep{claes2017decentralized} --- across six benchmark domains summarized in
Table~\ref{tab:benchmarks}.

\begin{table}[htbp]
\centering
\caption{Benchmark domains used for experimental evaluation.}
\label{tab:benchmarks}
\vspace{4pt}
\begin{tabular}{@{}p{3cm} p{2cm} p{4.5cm} p{2.5cm}@{}}
\toprule
\color{jhublue}\textbf{Domain} &
\color{jhublue}\textbf{Agents} &
\color{jhublue}\textbf{Description} &
\color{jhublue}\textbf{Metric} \\
\midrule
Search \& Rescue   & 6--12 & Locate targets in a grid with partial observability  & Targets found                \\
Formation Control  & 4--8  & Maintain a prescribed formation while avoiding obstacles& Formation error              \\
Cooperative Transport & 3--6 & Jointly move a heavy object to a goal location        & Transport time               \\
Area Coverage      & 8--20 & Maximize area swept by a sensor network               & \% area covered              \\
Pursuit-Evasion    & 4--10 & Capture adversarial targets with coordinated pursuit  & Capture rate                 \\
Warehouse Logistics& 5--15 & Route robots to fulfill pick-and-place tasks          & Tasks completed per minute   \\
\bottomrule
\end{tabular}
\vspace{4pt}
\end{table}

\section{Results}

Figure~\ref{fig:results} plots the key performance trends. BADC consistently
outperforms the decentralized baselines and approaches centralized performance
as the communication graph density increases beyond a modest threshold.

\begin{figure}[htbp]
\centering
\begin{tikzpicture}[scale=0.9]
  % Axes
  \draw[->, thick, color=jhuslate] (0,0) -- (10.5,0) node[right] {\footnotesize Communication rate (msg/agent/step)};
  \draw[->, thick, color=jhuslate] (0,0) -- (0,6.5) node[above] {\footnotesize Task completion rate};

  % Grid
  \foreach \y in {0,1,...,6} {
    \draw[color=jhulight, thin] (0,\y) -- (10,\y);
  }
  \foreach \x in {0,1,...,10} {
    \draw[color=jhulight, thin] (\x,0) -- (\x,6);
  }

  % Y-axis labels
  \foreach \y/\label in {0/0.0, 2/0.4, 4/0.8, 6/1.0} {
    \node[left,color=jhuslate,font=\footnotesize] at (0,\y) {\label};
  }

  % BADC (solid blue)
  \draw[line width=1.4pt, color=jhublue] plot coordinates {
    (0.5,0.55) (1.5,2.8) (2.5,4.2) (3.5,5.0) (5.0,5.5) (7.0,5.7) (10.0,5.9)
  };
  \node[color=jhublue, font=\footnotesize\bfseries] at (8.5,5.5) {BADC};

  % C-MDP (dashed gold, flat at top)
  \draw[line width=1.2pt, color=jhugold, dashed] plot coordinates {
    (0.5,2.0) (1.5,4.4) (2.5,5.3) (3.5,5.7) (5.0,5.85) (7.0,5.9) (10.0,5.95)
  };
  \node[color=jhugold, font=\footnotesize\bfseries] at (8.5,6.1) {C-MDP};

  % IQL (dotted gray)
  \draw[line width=1.0pt, color=gray, dotted] plot coordinates {
    (0.5,0.8) (1.5,1.0) (2.5,1.5) (3.5,2.2) (5.0,3.0) (7.0,3.8) (10.0,4.5)
  };
  \node[color=gray, font=\footnotesize] at (8.5,4.2) {IQL};

  % CBBA
  \draw[line width=1.0pt, color=gray, dashdotted] plot coordinates {
    (0.5,1.0) (1.5,1.8) (2.5,2.8) (3.5,3.5) (5.0,4.2) (7.0,4.7) (10.0,5.1)
  };
  \node[color=gray, font=\footnotesize] at (8.5,4.8) {CBBA};
\end{tikzpicture}
\caption{Task completion rate as a function of communication rate. BADC achieves
near-centralized performance with substantially less communication. Results
are averaged over 500 trials on the Search \& Rescue domain.}
\label{fig:results}
\end{figure}

\section{Physical Robot Experiments}

We deployed BADC on a fleet of eight TurtleBot3 Waffle Pi robots in a $120\,\text{m}^2$
warehouse environment at the Nimbus Institute for Robotics. The robots were
tasked with locating and transporting colored cubes to designated drop-off zones
under three communication regimes: full ($100\%$ connectivity), moderate ($60\%$
packet delivery), and degraded ($35\%$ packet delivery). Table~\ref{tab:physical}
summarizes the results.

\begin{table}[htbp]
\centering
\caption{Physical robot experiment results (mean $\pm$ std.\ over 30 trials).}
\label{tab:physical}
\vspace{4pt}
\begin{tabular}{@{}p{2.5cm} c c c@{}}
\toprule
\color{jhublue}\textbf{Regime} &
\color{jhublue}\textbf{Tasks Completed} &
\color{jhublue}\textbf{Time per Task (s)} &
\color{jhublue}\textbf{Collisions} \\
\midrule
Full ($100\%$)       & $18.4 \pm 1.2$  & $47.3 \pm 5.1$  & $2.1 \pm 1.3$  \\
Moderate ($60\%$)    & $17.1 \pm 1.8$  & $52.8 \pm 6.4$  & $3.4 \pm 1.8$  \\
Degraded ($35\%$)    & $14.6 \pm 2.3$  & $64.1 \pm 9.2$  & $5.2 \pm 2.6$  \\
\bottomrule
\end{tabular}
\vspace{4pt}
\end{table}

Even under severely degraded communication, BADC completed $79\%$ of the
full-connectivity task count, demonstrating the graceful degradation property
claimed in Chapter~\ref{ch:badc}.

% ======================================================================
% CHAPTER 6: CONCLUSION
% ======================================================================
\chapter{Conclusion}
\label{ch:conclusion}

\section{Summary of Contributions}

This dissertation introduced a principled framework for decentralized multi-robot
coordination that bridges the gap between the theoretical guarantees of centralized
planning and the practical robustness of distributed execution. The BADC framework
combines three ideas --- factored belief tracking, event-triggered consensus, and
submodular greedy planning --- into a single architecture with provable convergence
and approximation guarantees.

\section{Future Directions}

Several directions merit further investigation. \textbf{Heterogeneous teams},
in which robots possess different sensing and actuation capabilities, would
require extending the factored belief representation to capture capability-aware
intent models. \textbf{Adversarial settings}, including GPS spoofing and
communication jamming, raise questions about the security of consensus protocols
that this work has not addressed. \textbf{Learning from demonstration} offers a
promising avenue for reducing the planning overhead in repetitive tasks by
compiling frequently-used coordination patterns into reactive policies. Finally,
\textbf{human-in-the-loop coordination} --- where a human operator can inject
high-level guidance into an otherwise autonomous team --- poses interface and
authority-transfer challenges that blend robotics with human-computer interaction
research.

% ======================================================================
% APPENDICES
% ======================================================================
\appendix
\chapter{Convergence Proof Details}
\label{app:convergence}

\section{Preliminaries}

We restate several definitions from the theory of nonnegative matrices
\citep{seneta2006nonnegative} that are used in the proof of
Theorem~\ref{thm:convergence}.

\begin{definition}[Scrambling Matrix]
  A row-stochastic matrix $P \in \mathbb{R}^{n \times n}$ is \emph{scrambling} if,
  for every pair of rows $i, j$, there exists a column $k$ such that
  $P_{ik} > 0$ and $P_{jk} > 0$.
\end{definition}

\begin{definition}[Coefficient of Ergodicity]
  For a row-stochastic matrix $P$, the \emph{coefficient of ergodicity} is
  \begin{equation}
    \tau(P) = \frac{1}{2} \max_{i,j} \sum_{k=1}^n |P_{ik} - P_{jk}|.
  \end{equation}
  A matrix is scrambling if and only if $\tau(P) < 1$.
\end{definition}

\section{Main Argument}

The proof proceeds in three steps. First, we show that under the uniformly
jointly connected assumption, there exists a window size $B$ and a constant
$c < 1$ such that the product of any $B$ consecutive weight matrices has
coefficient of ergodicity at most $c$. Second, we apply the submultiplicativity
property of $\tau$ to bound the convergence rate. Third, we translate the
result in total-variation distance to the $\ell_1$ norm bound stated in
the theorem.

The complete derivation, including Lemmas A.1--A.5, is available in the
supplementary technical report \citep{calloway2026badc}.

% ======================================================================
% BIBLIOGRAPHY
% ======================================================================
\begin{thebibliography}{99}

\bibitem[Becker et~al.(2004)]{becker2004transition}
R.\ Becker, S.\ Zilberstein, V.\ Lesser, and C.\ V.\ Goldman.
\newblock Transition-independent decentralized Markov decision processes.
\newblock In \textit{Proc.\ AAMAS}, pages 41--48, 2004.

\bibitem[Bernstein et~al.(2002)]{bernstein2002complexity}
D.\ S.\ Bernstein, R.\ Givan, N.\ Immerman, and S.\ Zilberstein.
\newblock The complexity of decentralized control of Markov decision processes.
\newblock \textit{Math.\ Oper.\ Res.}, 27(4):819--840, 2002.

\bibitem[Calloway(2026)]{calloway2026badc}
M.\ E.\ Calloway.
\newblock Belief-augmented decentralized coordination for multi-robot systems.
\newblock Technical Report LASTR-2026-04, Laboratory for Autonomous Systems, 2026.

\bibitem[Choi et~al.(2009)]{choi2009consensus}
H.-L.\ Choi, L.\ Brunet, and J.\ P.\ How.
\newblock Consensus-based decentralized auctions for robust task allocation.
\newblock \textit{IEEE Trans.\ Robot.}, 25(4):912--926, 2009.

\bibitem[Claes et~al.(2017)]{claes2017decentralized}
D.\ Claes, F.\ Oliehoek, H.\ Baier, and K.\ Tuyls.
\newblock Decentralised online planning for multi-robot warehouse commissioning.
\newblock In \textit{Proc.\ AAMAS}, pages 492--500, 2017.

\bibitem[Foerster et~al.(2018)]{foerster2018counterfactual}
J.\ Foerster, G.\ Farquhar, T.\ Afouras, N.\ Nardelli, and S.\ Whiteson.
\newblock Counterfactual multi-agent policy gradients.
\newblock In \textit{Proc.\ AAAI}, 2018.

\bibitem[Hadjicostis and Dom{\'\i}nguez-Garc{\'\i}a(2017)]{hadjicostis2017average}
C.\ N.\ Hadjicostis and A.\ D.\ Dom{\'\i}nguez-Garc{\'\i}a.
\newblock Robust distributed average consensus via exchange of running sums.
\newblock \textit{IEEE Trans.\ Autom.\ Control}, 63(5):1492--1507, 2017.

\bibitem[Hansen et~al.(2004)]{hansen2004dynamic}
E.\ A.\ Hansen, D.\ S.\ Bernstein, and S.\ Zilberstein.
\newblock Dynamic programming for partially observable stochastic games.
\newblock In \textit{Proc.\ AAAI}, pages 709--715, 2004.

\bibitem[Kempe et~al.(2003)]{kempe2003gossip}
D.\ Kempe, A.\ Dobra, and J.\ Gehrke.
\newblock Gossip-based computation of aggregate information.
\newblock In \textit{Proc.\ FOCS}, pages 482--491, 2003.

\bibitem[Krause and Golovin(2014)]{krause2014submodularity}
A.\ Krause and D.\ Golovin.
\newblock Submodular function maximization.
\newblock In \textit{Tractability: Practical Approaches to Hard Problems},
  pages 71--104. Cambridge University Press, 2014.

\bibitem[Lowe et~al.(2017)]{lowe2017maddpg}
R.\ Lowe, Y.\ Wu, A.\ Tamar, J.\ Harb, P.\ Abbeel, and I.\ Mordatch.
\newblock Multi-agent actor-critic for mixed cooperative-competitive environments.
\newblock In \textit{Proc.\ NeurIPS}, pages 6379--6390, 2017.

\bibitem[Nair et~al.(2005)]{nair2005networked}
R.\ Nair, P.\ Varakantham, M.\ Tambe, and M.\ Yokoo.
\newblock Networked distributed POMDPs: A synthesis of distributed constraint
  optimization and POMDPs.
\newblock In \textit{Proc.\ AAAI}, pages 133--139, 2005.

\bibitem[Nemhauser et~al.(1978)]{nemhauser1978analysis}
G.\ L.\ Nemhauser, L.\ A.\ Wolsey, and M.\ L.\ Fisher.
\newblock An analysis of approximations for maximizing submodular set functions.
\newblock \textit{Math.\ Program.}, 14(1):265--294, 1978.

\bibitem[Olfati-Saber et~al.(2007)]{olfati2007consensus}
R.\ Olfati-Saber, J.\ A.\ Fax, and R.\ M.\ Murray.
\newblock Consensus and cooperation in networked multi-agent systems.
\newblock \textit{Proc.\ IEEE}, 95(1):215--233, 2007.

\bibitem[Ren and Beard(2005)]{ren2005consensus}
W.\ Ren and R.\ W.\ Beard.
\newblock Consensus seeking in multiagent systems under dynamically changing
  interaction topologies.
\newblock \textit{IEEE Trans.\ Autom.\ Control}, 50(5):655--661, 2005.

\bibitem[Seneta(2006)]{seneta2006nonnegative}
E.\ Seneta.
\newblock \textit{Non-negative Matrices and Markov Chains}.
\newblock Springer, 2nd edition, 2006.

\bibitem[Spaan and Vlassis(2005)]{spaan2005perseus}
M.\ T.\ J.\ Spaan and N.\ Vlassis.
\newblock Perseus: Randomized point-based value iteration for POMDPs.
\newblock \textit{J.\ Artif.\ Intell.\ Res.}, 24:195--220, 2005.

\bibitem[Tan(1993)]{tan1993multi}
M.\ Tan.
\newblock Multi-agent reinforcement learning: Independent vs.\ cooperative agents.
\newblock In \textit{Proc.\ ICML}, pages 330--337, 1993.

\end{thebibliography}

% ======================================================================
% BIOGRAPHICAL SKETCH (VITA)
% ======================================================================
\chapter*{Biographical Sketch}
\phantomsection
\addcontentsline{toc}{chapter}{Biographical Sketch}

Margaret Elena Calloway was born in Portland, Oregon, in 1994. She received
the B.S.\ degree in Mechanical Engineering with a minor in Computer Science
from the University of Washington in 2016, and the M.S.E.\ degree in Robotics
from the University of Pennsylvania in 2019. Her master's thesis, supervised
by Dr.\ Lina Thorvaldsen, investigated reinforcement learning for agile
quadruped locomotion on deformable terrain. She joined the doctoral program
in Mechanical Engineering at Johns Hopkins University in 2020 as a Vertex
Foundation Fellow. Her research interests include multi-agent systems,
distributed optimization, and field robotics. Following the completion of
her doctorate, she will join Apex Robotics as a Senior Research Scientist.

\end{document}
